\documentclass[11pt, a4paper]{article}

\usepackage[a4paper, top=2cm, bottom=2cm, left=2.5cm, right=2.5cm]{geometry}
\usepackage[T1]{fontenc}
\usepackage[utf8]{inputenc}
\usepackage{microtype}
\usepackage{parskip}
\usepackage{hyperref}

\hypersetup{
    colorlinks=true,
    linkcolor=blue,
    filecolor=magenta,      
    urlcolor=blue,
}

\begin{document}

\begin{flushleft}
\textbf{Quoc Anh Trinh (Andrew)} \\
Melbourne, Victoria \\
+61 451 886 468 \\
\href{mailto:andrews.trinh@gmail.com}{andrews.trinh@gmail.com} \\
\href{https://www.linkedin.com/in/qatrinh/}{linkedin.com/in/qatrinh} \\
\vspace{1em}
\today
\end{flushleft}

\vspace{1em}

\begin{flushleft}
\textbf{To the UpGuard Hiring Team,}
\end{flushleft}

\vspace{1em}

UpGuard’s mission to replace manual security bottlenecks with AI-driven precision resonates deeply with me. Processing 100 billion risk signals daily is an engineering feat that defines a category, and as an AI Engineer with a natural bias toward building and automating messy processes, I am excited to apply for the AI Engineer role.

Throughout my tenure at Bendigo \& Adelaide Bank, I have sat at the intersection of complex regulatory workflows and intelligent automation. I recently co-designed and delivered an end-to-end \textbf{AI-driven UAR-to-SMR pipeline} that leveraged \textbf{Pydantic-validated LLM prompting} and secure \textbf{data masking/rehydration} to reduce reporting preparation time by \textbf{80\%}. This project involved exactly the kind of operational deep dive described in the job post: mapping "As-Is" manual processes and translating them into a robust, high-integrity technical solution.

In my current role, I am leveraging AI to transform static regulation into active, high-performance governance. I use \textbf{Agentic AI} not just for information retrieval, but as specialized reasoning tools—such as the \textbf{typology extractor} I engineered to automatically transform 156 complex AUSTRAC documents into structured typologies. By integrating these agents into GCP-based data pipelines, I am replacing legacy manual bottlenecks with a transparent, automated framework that ensures 100\% continuity and logic alignment across millions of transactions.

My technical toolkit aligns directly with your stack:
\begin{itemize}
    \item \textbf{AI \& LLM Orchestration:} I have engineered Agentic AI solutions using LlamaIndex and FastAPI, including a \textbf{typology extractor} that transformed 156 guidance documents into structured monitoring data for framework comparison.
    \item \textbf{Automation \& Integration:} I built automated daily workflows that synchronized alerts across internal and vendor systems with \textbf{99.9\% data integrity}, eliminating manual transfers and ensuring seamless system orchestration.
    \item \textbf{Cloud \& Data Infrastructure:} I am comfortable working within \textbf{GCP} and have architected vector search infrastructure to power regulatory knowledge hubs, ensuring LLM outputs are grounded and actionable.
\end{itemize}

Beyond the technical skills, I thrive in the "discovery-to-deployment" phase. Whether it was leading the annual performance review logic for transaction monitoring or building typology management platforms, my focus has always been on turning ambiguous operational pain points into scalable, reliable tools that teams actually love to use.

UpGuard’s recent Series C highlights a scale and momentum that I find incredibly compelling. I am eager to bring my experience in building secure, transparent, and high-performance AI tools to help your Analytics team empower leaders with data-driven decision-making.

Thank you for your time and consideration. I look forward to the possibility of discussing how my background in automating complex banking workflows can translate into immediate impact for UpGuard.

\vspace{2em}

Sincerely,

\vspace{1em}

Quoc Anh Trinh (Andrew)

\end{document}
